%============================================================
% Chapter 28: Circular Arrays & Rotations
% STATUS: PLANNED  (generated by tools/outline.py; replace this file when the chapter is written)
%============================================================
\chapter{Circular Arrays \& Rotations}
\label{ch:circular}

\begin{tcolorbox}[enhanced, colback=white, colframe=secondaryblue, boxrule=1pt, arc=2mm]
\textbf{\large\textcolor{primaryblue}{An array whose end touches its start: doubling, the complement window, and rotations as shifts of the start index.}}
\end{tcolorbox}

\chapterfacts{1400}{2300}{5}{10}{30}

\begin{plannedbox}[Planned Contents --- this chapter is not written yet]
\textbf{Part A --- Lessons}
\begin{enumerate}[leftmargin=*, label=\textbf{A\arabic*.}, itemsep=1pt]
    \item Indexing modulo $n$ and the doubling trick
    \item Circular subarrays: total minus the complementary window
    \item Rotations: the best rotation and incremental rotation costs
    \item Circular two pointers and windows
    \item Cyclic shifts of permutations
\end{enumerate}
\textbf{Part B --- Patterns} (with candidate problems, checked when the chapter is written)
\begin{enumerate}[leftmargin=*, label=\textbf{P\arabic*.}, itemsep=1pt]
    \item \textbf{Doubling}
    \item \textbf{Complement window}
    \item \textbf{Best rotation}
    \item \textbf{Circular two pointers}
    \item \textbf{Cyclic shifts}
\end{enumerate}
\end{plannedbox}
